\documentclass[11pt,a4paper]{article}
\usepackage{notes}
\setlist{noitemsep,topsep=4pt}
\title{Week 1 Workshop — Gradients \& Curvature: From Topography to Hidden Fields}
\author{Geophysics IIIA (Potential Fields \& Geothermics)}
\date{\today}
\begin{document}
\maketitle

\section*{Purpose}
Build intuitive and mathematical comfort with gradients ($\nabla h$) and curvature ($\nabla^2 h$) in 2-D, starting from topography, then transferring the concepts to hidden subsurface scalar fields (e.g., density, temperature). This prepares students for potential fields (Weeks 3--7) and diffusion/heat flow (Weeks 8--10). \textit{Suggested pre-reading: Unit~1 notes on gradient, divergence, Laplacian, and Poisson/Laplace.} \emph{(See Unit1.pdf.)} \cite{unit1}


\section*{Learning outcomes}
By the end of the workshop, students will be able to:
\begin{enumerate}
  \item Draw and interpret the gradient vector field of a scalar map (magnitude and direction).
  \item Explain what the Laplacian measures (local curvature) and where it is large.
  \item Transfer gradient/curvature ideas from elevation to hidden fields; reason about subsurface contacts and sources/sinks.
  \item Articulate why moving farther from sources produces smoother maps (``distance acts like a low-pass filter'').
\end{enumerate}

\section*{Timing (180 min)}
\begin{description}
\item[0--10 min] \textbf{Concept warm-up.} Units of gradient; why a gradient is a vector (direction of steepest increase; magnitude = steepest slope).
\item[10--65 min] \textbf{Hands-on gradients (maps).} Using Mt.~Shasta and Marble Canyon topographic maps, students draw gradient arrows (longer where steeper), then compare to computed arrows/contours. Discuss where $|\nabla h|$ is large/small and why. \emph{(Based on gradient.pdf.)} \cite{gradient}
\item[65--90 min] \textbf{Mini-derivation \& vocabulary.} From slope to partial derivatives to $\nabla$; introduce $\nabla^2$ as divergence of $\nabla$ (curvature). Keep algebra light; focus on meanings and units. \emph{(See Unit1.pdf on grad/div/Laplacian.)} \cite{unit1}
\item[90--130 min] \textbf{Hidden fields page.} Treat the map as bulk density $\rho(x,y)$ (or temperature $T$). Sketch $\nabla \rho$ arrows and mark likely subsurface boundaries (faults, intrusive contacts). For each boundary, state whether the property is sharp or gradational and what that implies for $\nabla^2$ inside/outside sources.
\item[130--165 min] \textbf{Sharp vs. smooth.} Compare two maps: “close” vs “far” from a source. Identify visual cues (peak width, curvature, edge sharpness). Conclude: farther $\Rightarrow$ smoother; closer $\Rightarrow$ sharper.
\item[165--180 min] \textbf{Debrief.} How do these ideas appear in gravity/magnetics (continuation) and heat (diffusion)? Link to this week's practical (explicit Laplacian smoothing on ETOPO2) and to PS1's emphasis on units and assumptions. \emph{(See PS1.)} \cite{ps1}
\end{description}

\section*{Group structure}
Groups of 3--4 with rotating roles each week: Analyst (math/units), Modeler (figures), Interpreter (geologic meaning), Critic (assumptions/uncertainty).

\section*{Materials}
\begin{itemize}
  \item Printed A3/A4 maps from the gradient worksheet (Mt.~Shasta, Marble Canyon), blank overlays, fine markers. \cite{gradient}
  \item One-page “Hidden Fields” insert (below) for sketching density/temperature gradients and marking likely subsurface contacts.
\end{itemize}

\section*{Hidden Fields insert (student page)}
\textbf{1) Gradient of a hidden scalar.} Suppose the map shows shallow bulk density $\rho(x,y)$ instead of elevation. Draw arrows for $\nabla \rho$: longer where change is larger; direction toward increasing $\rho$. Label likely subsurface boundaries (faults, intrusive contacts, channel fills).

\textbf{2) Boundary reasoning.} For each boundary, is $\rho$ piecewise constant (sharp contact) or gradational? What sign/magnitude of $\nabla^2 \rho$ would you expect inside source regions vs outside (source-free) regions? \emph{(Recall from Unit1: outside sources $\nabla^2\Phi=0$; with sources $\nabla^2\Phi=-\rho/\kappa$ by analogy.)} \cite{unit1}

\textbf{3) Forward-model intuition.} Sketch equipotential contours $\Phi$ and field lines $\mathbf{F}=-\nabla\Phi$ for a simple dense plug beneath a flat surface. Where would a surface gravimeter see the largest vertical gradient?

\textbf{4) Heat-flow bridge.} If the map showed temperature $T$, which regions suggest internal heat production? Relate to the sign of $\nabla^2 T$.

\section*{Instructor notes}
Keep symbols/units consistent with the course notes; prioritize qualitative understanding. Emphasize that “farther $\Rightarrow$ smoother” returns in continuation (Weeks 6--7) and diffusion (Weeks 8--9). \cite{unit1}

\begin{thebibliography}{9}
\bibitem{unit1} Unit1.pdf: Course notes introducing work/potential, $\nabla$, $\nabla^2$, Poisson/Laplace. \cite{turn8search3}
\bibitem{gradient} gradient.pdf: Gradient worksheet with Mt. Shasta \& Marble Canyon maps and solutions. \cite{turn9search2}
\bibitem{ps1} geophysics\_ps1.pdf: Physical properties \& back-of-the-envelope problem set. \cite{turn8search2}
\end{thebibliography}
\end{document}